\input{../tex-inputs/latex-leseansicht-vorspann}

\section[Gustav Schwarzkopf an Arthur Schnitzler, 3. 5. 1897]{L04283 Gustav Schwarzkopf an Arthur Schnitzler, 3. 5. 1897}
\nopagebreak\mylabel{L04283v}
\rehead{ }\normalsize\beginnumbering\briefempfaengerindex{Schnitzler, Arthur@\textsc{Schnitzler, Arthur}!zzzSchwarzkopf, Gustav@\emph{von Gustav Schwarzkopf}!1897-05-031@{3. 5. 1897}|(be}
\toendnotes[C]{\smallbreak\pagebreak[2]}
\correspDesc{Versand  durch Gustav Schwarzkopf am 3. 5. 1897 in Wien
\newline{}Erhalt  durch Arthur Schnitzler im Zeitraum [4. 5. 1897 – 8. 5. 1897?] in Paris}\toendnotes[C]{\smallbreak}
\Standort{CUL, Schnitzler, B 96a.}
\physDesc{Briefkarte, 1330 Zeichen
\newline{}Handschrift: schwarze Tinte, deutsche Kurrent}
\pstart
           \raggedleft{}{\pb}3. V. 97\pend
           \vspace{0.5em}
\pstart
           Lieber Arthur! We{\geminationn} Sie vielleicht noch
               etwas Neid brauchen, um es deutlicher zu empfinden, daß es in Paris\oindex{Paris@\textbf{Paris}, \emph{Hauptstadt}|pw} augenblicklich angenehmer{ }ſein muß als in Wien\oindex{Wien@\textbf{Wien}, \emph{Verwaltungsgebiet}|pw},{ }ſo bin ich gerne bereit Sie zu beneiden. Gar
               zu{ }ſchwer wird mir das nicht werden und an Zeit fehlt es mir auch nicht. Von mir ka{\geminationn} ich Ihnen gar nichts mitteilen. \textsc{Zacconi\pwindex{Zacconi, Ermete 14.\,9.\,1857 Montecchio Emilia – 14.\,10.\,1948 Viareggio@\textsc{Zacconi, Ermete} (14.\,9.\,1857 Montecchio Emilia – 14.\,10.\,1948 Viareggio), \emph{Regisseur, Schauspieler}|pw}} iſt wirklich das einzig Bemerkenswerte, was ich{ }ſeit Ihrer Abreiſe erlebt habe.
               Hier iſt alles unverändert, nur \strikeout{\textcolor{gray}{ric}h} ein wenig öder und abgeſpielter, wie i{\geminationm}er am Schluß der Saiſon. Man führt noch i{\geminationm}er die{ }ſelben Geſpräche, man rechnet nich i{\geminationm}er den andern Leuten die \textsc{Tantièmen} nach, man{ }ſpielt noch i{\geminationm}er bei \textsc{Pucher\oindex{Wien@\textbf{Wien}!I., Innere Stadt@\textbf{I., Innere Stadt}!Café Pucher@\textbf{Café Pucher}, \emph{Kaffeehaus}|pw}}\textsc{Poker} und D\textsuperscript{r}\textsc{Eberma{\geminationn}\pwindex{Ebermann, Leo 16.\,7.\,1863 Draganovka – 9.\,10.\,1914 Wien@\textsc{Ebermann, Leo} (16.\,7.\,1863 Draganovka – 9.\,10.\,1914 Wien), \emph{Schriftsteller, Journalist, Rechtswissenschaftler}|pw}}{\pb}erkundigt{ }ſich noch i{\geminationm}er, ob man nicht bald nach Hauſe gehen wird. \textsc{Weiss\pwindex{Karlweis, Carl 23.\,11.\,1850 Wien – 27.\,10.\,1901 ebd.@\textsc{Karlweis, Carl} (23.\,11.\,1850 Wien – 27.\,10.\,1901 ebd.), \emph{Schriftsteller}|pw}} iſt Radfahrer geworden und beteiligt{ }ſich{ }ſelbſtverſtändlich{ }ſehr lebhaft an
               den intereſſanten Debatten über Räder, Geſchwindigkeit, Aufſteigen, Abſpringen u.ſ.w.
               Und ich ka{\geminationn} nicht mitreden! Meine Iſolirtheit gelangt
               mir i{\geminationm}er mehr zum Bewußtſein.\pend
           
\pstart
           Nun werden Sie hoffentlich nicht mehr bezweifeln, daß ich edel bin. Sie{ }ſchreiben mir
               9 Zeilen und meine Antwort wird die{ }ſchwindelnde Höhe von 27 Zeilen erreichen. Es
               bleibt mir auch nur noch eine Zeile für einen Gruß an D\textsuperscript{r}\textsc{Goldmann\pwindex{Goldmann, Paul 31.\,1.\,1865 Breslau – 25.\,9.\,1935 Wien@\textsc{Goldmann, Paul} (31.\,1.\,1865 Breslau – 25.\,9.\,1935 Wien), \emph{Schriftsteller, Journalist}|pw}} und für die Verſicherung, daß es mich{ }ſehr freuen wird Sie wieder in die Frankgaſſe\oindex{Wien@\textbf{Wien}!IX., Alsergrund@\textbf{IX., Alsergrund}!Frankgasse@\textbf{Frankgasse}, \emph{Straße}|pw} begleiten zu kö{\geminationn}en.\pend
           
\pstart
           Herzlichſt Ihr{\\[\baselineskip]}\spacefill\mbox{Gustav Schwarzk}\pend
           \leftskip=0em{}\selectlanguage{ngerman}\endnumbering\briefempfaengerindex{Schnitzler, Arthur@\textsc{Schnitzler, Arthur}!zzzSchwarzkopf, Gustav@\emph{von Gustav Schwarzkopf}!1897-05-031@{3. 5. 1897}|)be}\mylabel{L04283h}
\begin{anhang}
\end{anhang}\newcommand{\dateiname}{L04283}\newcommand{\titel}{Gustav Schwarzkopf an Arthur Schnitzler, 3. 5. 1897}\newcommand{\editorInnen}{Herausgegeben von Jahnke, SelmaMüller, Martin Anton}\input{../tex-inputs/latex-leseansicht-abspann}
